\section{De «se acabó» al cambio de proveedor de los clientes}

Esta sección entra en el punto más bajo de la empresa y en el punto de inflexión. En el video hay un segmento sobre «querer decir que se acabó»: el equipo tenía algo de tecnología láser, pero no un mercado lo bastante atractivo. Ese momento es decisivo, porque obliga al equipo a dejar de ser emprendedores tecnológicos y convertirse en verdaderos emprendedores de mercado: ya no se preguntan solo qué puede hacer la tecnología, sino por qué compra el cliente, cuándo compra y a qué precio.

\begin{importantbox}{El primer cambio de rumbo del emprendimiento en tecnología profunda}
Pasar de «tengo la tecnología» a «por qué el cliente tiene que comprar» es el cambio de rumbo decisivo del emprendimiento en tecnología profunda. Lo primero responde a una lógica de investigación y desarrollo; lo segundo, a una lógica de industria.
\end{importantbox}

\subsection{La lógica de precios: costo, valor y posición estratégica}

La subsección anterior planteó que el emprendimiento en tecnología profunda debe pasar de una lógica tecnológica a una lógica centrada en el cliente; esta subsección lleva ese cambio a la fijación de precios. El precio de un producto de hardware no se calcula como costo más margen. Debe considerar el costo de la BOM, el valor de seguridad para el cliente, el precio de los productos competidores que podrían sustituirlo, la reducción de costos esperada y la posición estratégica que da entrar en las armadoras. Si el precio inicial es incorrecto, se pueden perder clientes; si es demasiado bajo, quizá no alcance para sostener las entregas y la iteración.

\begin{figure}[H]
\centering
\includegraphics[width=0.96\textwidth]{pricing-logic.png}
\caption{La lógica de precios del hardware: el precio no es costo más margen, sino valor para el cliente, riesgo y posición en el mercado. Diagrama conceptual propio, elaborado a partir de la conversación de 01:20:47--01:38:15.}
\end{figure}

\begin{knowledgebox}{Lectura de la figura: el precio es un lenguaje estratégico}
El precio de un producto de hardware le comunica al cliente tres cosas a la vez: cuánto vale el producto, cuánto riesgo asume el proveedor y si en el futuro podrá producirse a escala. Si el precio inicial se calcula solo con base en el costo, se subestima el valor estratégico de entrar en la cadena de suministro de las armadoras.
\end{knowledgebox}

\begin{knowledgebox}{Las cinco variables del precio del hardware}
\begin{tabularx}{\textwidth}{p{0.20\textwidth}p{0.34\textwidth}X}
\toprule
Variable & Significado & Efecto \\
\midrule
Costo & Materiales, manufactura, rendimiento de producción, servicio posventa & Define el piso. \\
Valor & Mejora en seguridad, experiencia y desempeño & Define cuánto está dispuesto a pagar el cliente. \\
Competencia & Proveedores actuales y soluciones alternativas & Define el margen para entrar. \\
Escala & Reducción de costos con la producción en serie futura & Define la estrategia de precios a largo plazo. \\
Estrategia & Si se entra en la cadena de suministro de un cliente clave & Define si se puede sacrificar la utilidad a corto plazo. \\
\bottomrule
\end{tabularx}
\end{knowledgebox}

\begin{warningbox}{La trampa de los precios bajos en hardware}
Un precio bajo puede acelerar que los clientes cambien de proveedor, pero también puede dejar a la empresa con margen bruto negativo, carga de servicio posventa y presión sobre la calidad durante mucho tiempo. El precio del hardware debe servir a la vez para entrar al mercado y para mantener viva a la organización; no puede perseguir solo «ganar al cliente».
\end{warningbox}

\subsection{Los clientes empiezan a cambiar de proveedor}

El precio muestra cómo el proveedor expresa su valor; esta subsección examina cómo votan de verdad los clientes. Que un cliente deje a su proveedor actual por uno nuevo es el momento de ruptura en la industria del hardware. No se logra con una sola demostración, sino con la acción conjunta del costo, el desempeño, la confiabilidad, la capacidad de entrega y las señales de la industria. Que el cliente cambie de proveedor significa que cree que la nueva solución no solo es más barata, sino suficientemente confiable.

\begin{figure}[H]
\centering
\includegraphics[width=0.96\textwidth]{customer-defection.png}
\caption{Los clientes empiezan a cambiar de proveedor: el momento de ruptura en la cadena de suministro de hardware llega cuando el cliente está dispuesto a dejar la solución anterior por una nueva. Diagrama conceptual propio, elaborado a partir de la conversación de 01:38:15--01:58:07.}
\end{figure}

\begin{knowledgebox}{Por qué los clientes cambian de proveedor}
\begin{tabularx}{\textwidth}{p{0.20\textwidth}p{0.34\textwidth}X}
\toprule
Condición & Qué ve el cliente & Qué significa para el proveedor \\
\midrule
Desempeño confiable & Cumple con los escenarios críticos & La tecnología ya no es solo una demo. \\
Ventaja de costo & El costo total baja de forma notable & Tiene la oportunidad de entrar en modelos de producción en serie. \\
Entrega estable & Puede suministrar según lo programado & Se valida la capacidad de la organización. \\
Riesgo de sustitución controlable & Migrar desde la solución anterior no se sale de control & El cliente está dispuesto a asumir el riesgo del cambio. \\
\bottomrule
\end{tabularx}
\end{knowledgebox}

\subsection{Resumen de la sección}

Para pasar de «se acabó» a que los clientes cambien de proveedor, lo decisivo es pasar de la confianza en la tecnología propia a la validación por parte del cliente. El verdadero punto de inflexión de una empresa de tecnología profunda llega cuando el cliente está dispuesto a confiarte su riesgo.
